\Needspace{12\baselineskip}
\SourceLeaf{93}%

\begin{center}
BEMERKUNGEN ZUR »ENTWICKLUNG VON $(h-\cos\varphi)^{-\frac{3}{2}}$
IN DIE REIHE $A^{(0)}+2A^{(1)}\cos\varphi+2A^{(2)}\cos2\varphi+\cdots$«.
\end{center}

Diese Entwicklung von $(h-\cos\varphi)^{-\frac{3}{2}}$ stammt vermuthlich aus der Zeit um
1813. Das besondere Interesse derselben ist in ihren mehrfachen Beziehungen zu sonstigen
Untersuchungen von \textsc{Gauss} begründet.

Für die beiden ersten Coefficienten $A^{(0)}$, $A^{(1)}$ hat man
zunächst die folgenden Integraldarstellungen:

\[
A^{(0)}=\int_{0}^{2\pi}\frac{d\varphi}{2\pi(h-\cos\varphi)^{\frac{3}{2}}},\qquad
A^{(1)}=\int_{0}^{2\pi}\frac{\cos\varphi\,d\varphi}{2\pi(h-\cos\varphi)^{\frac{3}{2}}}.
\]

\noindent Substituirt man hier $\varphi=2T$ und setzt $h-1=m^{2}$, $h+1=n^{2}$, so ergiebt sich:

\[
A^{(0)}=\int_{0}^{2\pi}\frac{dT}{2\pi(mm\cos T^{2}+nn\sin T^{2})^{\frac{3}{2}}},
\]

\[
A^{(1)}=\int_{0}^{2\pi}\frac{(\cos^{2}T-\sin^{2}T)\,dT}{2\pi(mm\cos T^{2}+nn\sin T^{2})^{\frac{3}{2}}}.
\]

\noindent Die Berechnung von $A^{(0)}$ und $A^{(1)}$ aus $m$ und $n$ durch den
Algorithmus des arithmetisch-geometrischen Mittels hat \textsc{Gauss} in den Artikeln 16 ff.\
seiner Abhandlung »Determinatio attractionis etc.« (ges. Werke III p.\,331 ff.)
geleistet. Die daselbst in Art.\,19 mit $P$ und $Q$ bezeichneten Integrale liefern:

\[
A^{(0)}=P+Q,\qquad A^{(1)}=P-Q.
\]

\Needspace{10\baselineskip}
\noindent Um die damaligen mit den hier vorliegenden (bei Berechnung des
arithmetisch-geometrischen Mittels benutzten) Bezeichnungen in Einklang zu setzen, hat
man an $m^{2}=h-1$, $n^{2}=h+1$ anzuknüpfen und findet in $t$:

\[
m=\frac{1-t}{\sqrt{2t}},\qquad n=\frac{1+t}{\sqrt{2t}}.
\]

\noindent Die Zahlen $a$, $b$, von denen hier das Mittel $a^{(\infty)}$ gebildet ist, sind
demnach:

\[
a=\Bigl(\frac{1-t^{2}}{2t}\Bigr)^{2}m',\qquad b=\Bigl(\frac{1-t^{2}}{2t}\Bigr)^{2}n',
\]

\noindent und es gilt:

\[
a^{(\infty)}=\Bigl(\frac{1-t^{2}}{2t}\Bigr)^{2}\mu,
\]

\noindent wo $\mu$, wie a.\,a.\,O., das arithmetisch-geometrische Mittel zwischen $m$ und $n$ ist.
Die am Schlusse von Art.\,17 der genannten Abhandlung angegebene Gleichung für $\nu$
schreibt sich auf Grund der ebenda nächst voraufgehenden Relationen:

\[
\nu=2\frac{\lambda'}{m'}\Bigl(1+2\frac{\lambda''}{m''}\Bigl(1+2\frac{\lambda'''}{m'''}(1+\cdots\Bigr\}.
\]
